\documentclass[12pt]{article}
\usepackage[margin=1.0in]{geometry}
\usepackage{setspace}
\setstretch{2}
\usepackage{graphicx}
\usepackage{algorithm}
\usepackage[noend]{algpseudocodex}
\algnewcommand\algorithmicforeach{\textbf{for each}}

\usepackage{soul}
\soulregister\citep7
\soulregister\citeauthor7

\usepackage{hyperref}
\hypersetup{%
    colorlinks=true,
    linkcolor=black,
    filecolor=black,
    urlcolor=black,
    citecolor=blue,
}
\usepackage{siunitx}
\DeclareSIUnit{\molar}{M}

% Bibliography
\usepackage{natbib}
\bibliographystyle{msb3} % the optional note field in the bib file is useful for MSB data refs

\begin{document}

\begin{titlepage}
    \begin{center}
        \vspace*{1cm}

        {\Large Atomic elementary flux modes explain the steady state flow of
        metabolites in large-scale flux networks}

%        \vspace{1.0cm}
%        aEFMs explain metabolite flows
            
        \vspace{1.0cm}

        Justin G. Chitpin$^{1,2}$ and Theodore J. Perkins$^{1,2*}$

        \vspace{1.0cm}

        $^1$Ottawa Hospital Research Institute, 501 Smyth Road, Ottawa, K1H 8L6, Ontario, Canada\\
        $^2$Ottawa Institute of Systems Biology, Department of Biochemistry, Microbiology and Immunology, University of Ottawa, 451 Smyth Road, Ottawa, K1H 8M5, Ontario, Canada
           
        \vspace{0.5cm}

        $^*$ Corresponding author.

        \vspace{1.0cm}
     
        Flux network/Elementary flux mode/flow decomposition/Markov chain/atom mapping
            
   \end{center}
\end{titlepage}

\clearpage

\section*{Abstract} % 175 words max

Steady state fluxes are a measure of cellular activity under metabolic
homoeostasis, but understanding how individual substrates are metabolised
remains a challenge in large-scale networks. Pathway-based approaches such
as elementary flux mode (EFM) analysis are limited to small networks due
to the combinatorial explosion of pathways and the ambiguity of
decomposing fluxes onto EFMs. Here, we present an alternative approach to
explain metabolic fluxes in terms of the steady state flow of their atomic
constituents. We refer to these pathways as atomic elementary flux modes
(AEFMs) and show that computations involving AEFMs are orders of magnitude
faster than standard EFMs. Using our approach, we enumerate carbon and
nitrogen AEFMs in five genome-scale metabolic models and compute the AEFM
decomposition of fluxes estimated in a HepG2 liver cancer cell line. Our
results systematically characterise carbon and nitrogen remodelling and,
on the HepG2 network, predict glutamine metabolism through a recently
discovered non-canonical TCA cycle.

\clearpage

\section*{Introduction} % intros tend to be very succinct

% (a) flux overview aasdf
Recent technological and computational advancements are leading to a rise in
metabolic flux
data~\citep{SternAlon2023Imac,LeeWonDong2019Spas,WangLin2022Srit,WangRuohong2022Gstm}.
These fluxes quantify the rate of metabolite interconversions within
a cell or their movements across subcellular compartments. Metabolic
fluxes are studied to understand the dynamics of substrate
switching~\citep{ZampieriMattia2012Rmuc,KochanowskiKarl2021Gcom}, redox
balancing~\citep{McKinlayJamesB2010Cdfa,FranchinaDavideG2022Grbc},
overflow
metabolism~\citep{BasanMarkus2015OmiE,VanderHeidenMatthewG2009UtWE}, cell
differentiation~\citep{EndoYusuke2022Fami,BenjaminJacksonT2024Mrot}, and
cellular
communities~\citep{SchuetzRobert2012MOoM,YuJasonSL2022Mcfr,Gustafsson2024Mcbc}.
Hence, `fluxomics'~\citep{EmwasAbdul-Hamid2022F-NM,SanfordKarl2002Gtfa} is
emerging as a field dedicated to quantifying the totality of fluxes within
biological systems.

% (b) estimating steady state fluxes - happy with the writing
Current experimental and computational methods generally estimate fluxes
under metabolic homeostasis. Under this assumption, the production and
consumption of metabolites are balanced over time, and the fluxes are said
to be under steady state. Techniques such stable isotope tracing analysis
(SITA) experimentally determine steady state fluxes by culturing cells in
labelled substrates and quantifying their incorporation within endogenous
metabolites~\citep{BuescherJoergM2015Arfi,AntoniewiczMaciekR2018Agt1}.
Using this method, fluxes have been measured in major metabolic subsystems
such as glycolysis, pentose phosphate pathway, and the TCA
cycle~\citep{ZamboniNicola20091mfa,AhnEunyong2017Tfro,WuChao2019FARC}. Many
computational methods, including but not limited to flux balance analysis
(FBA), have also been proposed to infer steady state fluxes from transcript,
protein, and metabolite abundance
data~\citep{UematsuSaori2022Mlmf,GonzalezArrueNicolas2023Peom,HuangYuefan2023Ccmf,KasteJoshuaAM2023Afpu,LeeJustinY2023Tiac}.
Collectively, these approaches are being used to estimate and predict fluxes in
genome-scale metabolic models (GEMs) involving hundreds to thousands of
reactions~\citep{NilssonAvlant2020Qaoa,CherkaouiSarah2022Afao,GelbachPatrickE2024Fsig}.

% (c) introducing problem
Although much effort has been made to estimate steady state fluxes, less
emphasis has been placed on developing methods to analyze these
large-scale flux datasets. In contrast to abundance data (e.g.\
transcripts, proteins, metabolites), steady state fluxes are constrained
by the reaction stoichiometries between metabolite inflows and outflows.
In a network with fixed source and sink fluxes, an increase in flux along
one metabolic pathway must decrease fluxes along one or more other
pathways to maintain mass conservation. For simple networks, one may
predict how the system would evolve by identifying the up and downstream
reactions flanking an enzymatic perturbation. For example,
\citeauthor{SchwartzJean-Marc2006Qema} were able to
enumerate all glycolytic pathway fluxes in a 12-reaction model of yeast
glycolysis~\citep{SchwartzJean-Marc2006Qema}. Yet in more complex genome-scale
metabolic networks containing thousands of fluxes
(e.g.\ \cite{BrunkElizabeth2018Reat,RobinsonJonathanL2020Aaoh,WangHao2021Gmnr}),
there exists no hierarchical relationship between reactions. A change in one
reaction rate, let alone multiple perturbations, could alter network fluxes well
beyond neighbouring reactions~\citep{NobileMarcoS2021Agsa}.

% (d) EFMs
This problem of understanding the flow of metabolites within steady state
flux networks has historically been addressed through the concept of
elementary flux modes (EFMs). EFMs are a pathway-based approach to analyze
metabolic networks and are defined as minimal sets of biochemical
reactions that maintain steady state flux~\citep{SchusterStefan1994Oefm}.
This minimal property specifies that an EFM cannot be decomposed into two
or more smaller pathways carrying steady state flux. An attractive
property of EFMs is that any set of steady state fluxes in a network can
be explained as a positive, linear combination of
EFMs~\citep{SchusterStefan1999Doef}. While EFMs are a mechanistic
explanation of how metabolites travel through networks, the number of EFMs
scales combinatorially as a function of network
size~\citep{AcunaVicente2009Maci}. It thus remains computationally
infeasible to enumerate EFMs in GEMs after over a decade of algorithmic
advancements (e.g.\
\citealt{GagneurJulien2004Coem,TerzerMarco2008Lcoe,BuchnerBiancaA2021EaPp}).
The number of EFMs also poses a challenge for explaining network fluxes in
terms of their EFMs. When there are more EFMs than linearly independent
fluxes, there exist multiple ways to assign EFM weights that reconstruct
the observed network fluxes. Consequently, structural analyses of EFMs and
downstream analyses of their weights have been limited to small-scale
metabolic
subnetworks~\citep{FakihIbrahim2023Dgmm,MpabanziLiliane2022Frca,ToyaYoshihiro2022Mpef,MatsudaFumio2011Esoy,SchwartzJean-Marc2006Qema}.

% (e) Network notions
In parallel with EFM analysis, the task of identifying metabolic
pathways has also been approached by numerous other graph-theoretic
approaches. These methods derive graphs based, in part, on the metabolic
flux network (e.g.\ Figure~\ref{fig01:Notions}a) that highlight specific
relationships between substrates and products. The simplest
representation, for example, is the metabolic
graph~\citep{AritaMasanori2000Mrus,WagnerAndreas2001Tswi}
(Figure~\ref{fig01:Notions}b), whereby metabolite nodes are connected to
each other based on substrate-product reaction relationships. These graphs
are typically analysed by identifying the $K$-shortest pathways between
source and target metabolites. However, these methods often returned
biologically meaningless pathways due to metabolite byproducts and
co-factors, leading to pathways that inaccurately described the
remodelling of source/target metabolite pairs. Subsequent
refinements implemented heuristics to avoid enumerating these
pathways~\citep{McShanD2003PathMiner,RahmanSA2004Mpaw,CroesDidier2006IMPi,PeyJon2011Pfma,HafnerJasmin2021NFmp};
however, this problem was fully solved by the introduction of atomic
transition graphs~\citep{AritaMasanori2003ISAT}. These graphs use
atom mapping data to model the movement of individual atoms across
substrates and products, thereby avoiding metabolic byproducts
(Figure~\ref{fig01:Notions}c). We refer to weighted versions of these
graphs as \emph{atomic flux graphs} which are used in isotope-based
metabolic flux analysis to predict intracellular reaction
fluxes~\citep{WiechertWolfgang1999Brsi,WeitzelMichael2007Ttom}.

%However due to the number of potential paths in these metabolite-atom
%graphs, current methods only compute $K$-shortest pathways which may not
%represent 

%methods to analyse these graphs rely on $K$-shortest path routing
%algorithms to identify pathways, potentially missing out on longer
%biologically relevant pathways.

%Methods to analyse these graphs can now compute simple cycles and
%branching paths involving source/target metabolites; however, these
%pathways are still limited by $K$-shortest path routing
%algorithms~\citep{BlumTorsten2008Mfsf,PitkanenEsa2009Ibpi,HeathAllisonP.2010Fmpu,HeathAllisonP.2011AAfE,PeyJon2011Pfma,LatendresseMario2014Omrs,KimSarahM.2020Itoa}.

%Many refinements have since been made to identify $K$-shortest simple
%paths/cycles and branching paths involving source/target metabolite
%pairs~\citep{BlumTorsten2008Mfsf,PitkanenEsa2009Ibpi,HeathAllisonP.2010Fmpu,HeathAllisonP.2011AAfE,PeyJon2011Pfma,LatendresseMario2014Omrs,KimSarahM.2020Itoa}.
%When reaction fluxes are known, the corresponding atomic fluxes can be
%assigned to the directed edges, converting them into atomic flux graphs
%used in isotope-based metabolic flux analysis to predict intracellular
%reaction fluxes~\citep{WiechertWolfgang1999Brsi,WeitzelMichael2007Ttom}.

% (f) EFM/AEFM segue with Pey's work
Current methods to analyse metabolic flux networks and atomic flux graphs
are often viewed as independent approaches for metabolic pathway
analysis~\citep{deFigueiredoLuisF.2009Csbp,AlgfoorZeyadAbd2017Iomp}.
However, both can be viewed as related types of flux networks operating at
the metabolite or metabolite-atom level. One could envision applying
metabolic pathway analysis techniques designed for metabolic flux networks
on their corresponding atomic flux graphs. This ideas was previously
explored by \citeauthor{PeyJon2011Defm}~\citep{PeyJon2011Defm} who proposed
the concept of elementary carbon modes (ECMs). These ECMs were defined as
minimal, mass-balanced pathways tracing the flow of individual carbon
atoms in an atomic flux graph~\citep{PeyJon2011Defm,PeyJon2012Itmd}. Like
paths in metabolite-atom graphs, ECMs correspond to either simple paths
spanning source and sink metabolite carbons or simple cycles involving
internal metabolite-carbon positions.

%EFMs, a positive linear combination of ECMs is aThey also showed that the
%totality of ECM weights was guaranteed to reh

%Because ECMs are EFMs defined with respect to carbon flux networks, the
%totality of ECM weights is guaranteed to reconstruct the overall network
%carbon mass flow~\citep{PeyJon2011Defm}.

% (g) AEFM definition plus EFM/AEFM toy example
Since ECMs are specific to carbon pathways, we propose generalising ECMs
to any element in an atomic flux network. We refer to these pathways as
\emph{atomic elementary flux modes} (AEFMs) and formally define them as
the EFMs of an atomic flux graph. To avoid confusion with classical EFMs
of metabolic flux networks, we henceforth refer to these pathways as
\emph{molecular} EFMs to emphasize their balance of molecular
stoichiometries. Figure~\ref{fig01:Notions}d shows the molecular EFMs for
the metabolic flux network in Figure~\ref{fig01:Notions}a, while
Figure~\ref{fig01:Notions}e shows the AEFMs for the corresponding atomic
transition/flux graph. Although there are a different number of molecular
EFMs-versus-AEFMs, both pathways carry some form of steady state flux. The
molecular EFMs balance the reaction fluxes and molecular stoichiometries,
while the AEFMs balance the atomic fluxes and atom mass transfers. We
further observe that the flux through each reaction is equal to the
underlying fluxes at the atomic level~\citep{PeyJon2012Itmd}. This leads
to the following linear relationship between molecular EFM and AEFM
weights: $w_1 = x_1 = x_2$, and $w_2 = x_3 = x_4 = y_1$. Thus in
Figure~\ref{fig01:Notions}d-e, knowledge of either $x_1$/$x_2$ and
$x_3$/$x_4$/$y_1$ in the atomic flux graph is sufficient to determine the
EFM weights in the original metabolic flux network.

% (h) Linker paragraph
Although corresponding AEFM and molecular EFM weights presents a new
computational challenge, we focus our attention on two fundamental
problems associated with AEFMs. Following molecular EFMs, these problems
are enumerating and uniquely decomposing atomic fluxes onto AEFMs.
Addressing the latter problem is relevant since the number of AEFMs can
exceed the number of linearly independent fluxes in more complex atomic
flux graphs. Indeed, \citeauthor{PeyJon2012Itmd} encountered this latter
problem with ECMs and were unable to uniquely estimate their weights
beyond simple carbon flux graphs~\citep{PeyJon2012Itmd}.

% (i) Reference to CHMC
In our previous work, we developed a novel solution to the EFM flux
decomposition problem for the special case of unimolecular reaction
networks under a Markov constraint. We modelled a hypothetical particle
moving between metabolites under the assumption that the transition
probability of each reaction was proportional to the observed steady state
flux. By expanding the notion of Markov chain state to include sequences
of metabolites, we proposed an algorithm to efficiently compute simple
cycle probabilities corresponding to a given EFM. We called this construct
the \emph{cycle-history Markov chain} (CHMC) and proved that these EFM
probabilities could be scaled to weights that fully reconstructed the flux
along each reaction~\citep{ChitpinJustinG2023AMct}.

% (j) directives
Here, we combine three ideas in metabolic modelling to produce a pipeline
for analysing metabolic flux networks under steady state. First, we use
the state-of-the art atom mapping program RXNMapper (assisted by manual
curation) to automatically construct atomic flux graphs for any given
metabolic network. Second, we merge the concept of metabolite-atom paths
with elementary flux modes through our proposal of AEFMs. Third, we expand
the CHMC method to operate on atomic transition/flux graphs to enumerate
and compute unique AEFM weights. Using this pipeline, we demonstrate the
feasibility of AEFM enumeration by showing that the number of these carbon
and nitrogen pathways are orders of magnitude smaller than the number of
molecular EFMs in five metabolic networks. We then analyse the structure
of these AEFMs characterising carbon and nitrogen remodelling within and
across nutrient sources. Lastly, we run our full pipeline on a HepG2
metabolic flux network to examine glutamine carbon remodelling. Our AEFM
weight analysis reveals that glutamine consumption is routed through
a small number of highly active AEFMs, including a recently discovered
non-canonical TCA cycle~\citep{ArnoldPaigeK2022Anta}. 

\section*{Results}

\subsection*{Pipeline for AEFM enumeration and weight identification} % pipeline overview


% Describe pipeline but now on more realistic networks.

% Clarification needed on the atomic flux _subgraph_ bit

% Can mention the atom never splitting bit

We first describe what molecular EFMs look like in biologically-relevant
metabolic networks and provide intuition why enumerating and decomposing
fluxes onto these pathways remains a fundamental problem. Figure
\ref{fig02:Pipeline}a shows an example molecular EFM involving substrates
glucose and ATP which are consumed to produce metabolites G3P, DHAP, and
ADP. Multiple source and sink fluxes are required to balance the internal
metabolite stoichiometries, leading to several branching and converging
paths within the molecular EFM. As the metabolic networks expands to
include dozens or more reactions, the corresponding molecular EFMs can
grow increasingly complex with numerous source and sink fluxes required to
stoichiometrically balance participating metabolites. In particular, there
may be multiple sets of reactions involved to balance pairs of cofactors
(e.g.\ ATP/ADP, NAD$^+$/NADH), metabolic byproducts (e.g.\ water, hydrogen
ions), or coenzymes (e.g.\ coenzyme A). This complexity can lead to an
intractable number of pathways that may only differ by reaction subsets
required to mass-balance these ubiquitous metabolites. Removing these
metabolites may, in part, solve this problem; however, the choice of
metabolite to prune is often subjective and not guaranteed to sufficiently
reduce the total number of molecular EFMs.

Given many metabolic flux networks are generated to study nutrient
catabolism~\citep{ZamboniNicola2015DtMS} or engineer desirable metabolic
products~\citep{ZanghelliniJurgen2013Efmi}, we argue for studying source
metabolite-derived AEFMs to reduce the combinatorial number of molecular
EFMs. We highlight an example carbon AEFM in Figure \ref{fig02:Pipeline}b
which shows how the C6 glucose carbon is remodelled during the first steps
of glycolysis. As atoms are never split apart or fused together in
metabolic reactions, these AEFMs always correspond to simple paths
spanning source and sink metabolites or simple cycles between internal
metabolites. These pathways are straightforward to interpret because each
atom within an AEFM can only transit from a single substrate to product,
regardless of the network size, connectivity, or number of participating
substrates/products in each reaction. Hence, this property enables us to
enumerate and compute the AEFM weights using our previously described CHMC
method to uniquely decompose steady state fluxes onto EFMs for
unimolecular reaction networks~\citep{ChitpinJustinG2023AMct}.

We outline our computational pipeline in Figure \ref{fig02:Pipeline}c to
enumerate AEFMs and uniquely assign their weights when the steady state
network fluxes are known. In step 1, we require as inputs the metabolic
network defined by a set of metabolites and reactions encoded as
a stoichiometry matrix, the metabolite structures denoted as SMILES
strings, and the steady state fluxes. From the stoichiometry matrix and
SMILES strings, we construct reaction SMILES strings which are imputed
into the atom mapping program RXNMapper~\citep{Schwaller2021extraction} to
map the position of each atom within each reaction equation.

%\hl{

Using these atom mapping predictions, we construct atomic transition/flux
subgraphs rooted on each source metabolite atom entering the metabolic
network. Figure~\ref{fig02:Pipeline}c shows an example atomic
transition/flux graph for the C6 carbon of input glucose entering the
network with atomic fluxes proportional to the molecular fluxes,
accounting for reactions with multiple substrate stoichiometries
(see Methods and Figures~{\ref{figS01:Rules1}}-{\ref{figS02:Rules2}}). In
step 3, we construct the corresponding ACHMC to enumerate the AEFMs for
each atomic transition/flux graph. Finally, the AEFM weights are computed
in step 4 by steady state analysis of the ACHMC. 
%Further pipeline details are described in the Methods.
%}

\iffalse
While it would be logical to construct atomic transition/flux graphs, as
shown in Figure~{\ref{fig01:notions}}c, we use the atom mapping
predictions to extract the relevant metabolite-atom positions transited by
each source metabolite atom entering the metabolic network. This allows us
to directly construct the atomic transition subgraphs with transi

Instead of constructing an atomic flux graph shown in
Figure~{\ref{fig01:notions}}c, we use the atom mapping data to extract the
relevant metabolite-atom positions transited by each source metabolite
atom entering the metabolic network. This allows us to construct multiple,
albeit smaller, atomic flux subgraphs with transition probabilities
proportional to the atomic fluxes through the underlying reactions,
accounting for reactions with multiple substrate stoichiometries
(see Methods and Figures~{\ref{figS01:Rules1}}-{\ref{figS02:Rules2}}).
%For example, 10\SI{}{\milli\molar} of input glucose flux would correspond
%with 10\SI{}{\milli\molar} of atomic flux through each glucose carbon.
\fi

\subsection*{Enumerating AEFMs is computationally tractable in large-scale networks}

Given the major computational challenge of enumerating molecular EFMs in
networks with $>100$ metabolites and reactions, we first sought to test
whether it was feasible to enumerate AEFMs in GEMs. We selected five GEMs
across distinct organisms in ascending order of network size. These
networks include a metabolic model of \textit{E. coli} (E. coli core;
\citealt{OrthJeffreyD2010RaUo}; Data ref:
\citealt{OrthJeffreyD2010RaUoDATASET}), human red blood cells (iAB RBC
283; \citealt{BordbarAarash2011iApd}; Data ref:
\citealt{BordbarAarash2011iApdDATASET}), \textit{H. pylori} (iIT341;
\citealt{ThieleInes2005EMRo}; Data ref:
\citealt{ThieleInes2005EMRoDATASET}), \textit{S. aureus} (iSB619;
\citealt{BeckerScottA2005Grot}; Data ref:
\citealt{BeckerScottA2005GrotDATASET}), and human liver cancer cells
(HepG2; \citealt{NilssonAvlant2020Qaoa}; Data ref:
\citealt{NilssonAvlant2020QaoaDATASET}). These models were pre-processed
to ensure unambiguous atom mappings across all reactions. Namely, all
metabolites with no known structures (pseudometabolites), and
pseudoreactions with non-integer stoichiometries were removed from the
networks. The resulting metabolic networks ranged between 71-547
metabolites and 144-974 reactions (Figure \ref{figS03:Hairballs}). We then
attempted to enumerate the molecular EFMs using FluxModeCalculator, which
is the fastest molecular EFM enumeration
program~\citep{vanKlinkenJanBert2016Faet}, and our AEFM enumeration
pipeline (Figure~\ref{fig01:Notions}c steps 1-3). We first compiled the
metabolite SMILES strings and constructed ACHMCs rooted on each carbon and
nitrogen atom across all source metabolites. The resulting summary
statistics describing the GEMs and ACHMCs are presented in Table 1. Across
all carbon ACHMCs, we observed a broad increase in the total number of
AEFMs as a function of metabolic network size. The one exception to this
observation was the HepG2 network which contained fewer metabolites and
reactions than iSB619, yet exhibited two orders of magnitude more carbon
AEFMs. Inspection of the HepG2 network revealed it only contained 40
source metabolites compared to the 198 source metabolites in the iSB619
network, indicating a greater degree of network connectivity between
internal metabolites (Figure \ref{figS03:Hairballs}).

Figure \ref{fig03:Benchmark}a shows the total number of AEFMs and
molecular EFMs enumerated by our program MarkovWeightedEFMs.jl and
FluxModeCalculator on the five GEMs. Using our ACHMC pipeline, we
enumerated between $10^3 - 10^6$ carbon and $10^2-10^4$ nitrogen AEFMs
within each GEM. In contrast, FluxModeCalculator failed to enumerate
molecular EFMs in all but the two smallest E. coli core and iAB RBC 283
datasets, returning $10^3$ times more molecular-versus-atomic EFMs. 

By focusing on an atom of interest, our ACHMC approach greatly reduced the
number of EFMs to biologically relevant pathways involving nutrient source
remodelling. This reduction in pathways resulted in AEFM-versus-EFM
enumeration completing $10^3-10^4$ times faster (Figure
\ref{fig03:Benchmark}b). While running FluxModeCalculator in parallel did
improve run times by 8.2-8.6 times (Figure \ref{figS04:Scaling}), it
remained over $10^1-10^2$ slower than enumerating the corresponding carbon
and nitrogen AEFMs. One may think that expanding metabolites into
metabolite-atom positions would increase the number of ACHMC states and
therefore the computational run time. However, we found that AEFM
enumeration actually involved fewer metabolite-atom positions than the
total number of metabolites in a given network. Constrained by the atom
mapping predictions, we observed that the average carbon and nitrogen
ACHMC across all five GEMs only transited
0.41\% and 1.7\% of the total number of metabolite-carbon and
metabolite-nitrogen positions, respectively (Figure
\ref{figS05:StateSpace}). While serial run times in Figure
\ref{fig03:Benchmark}c scaled quadratically (Figure
\ref{fig03:Benchmark}c; $\text{slope}=1.85, \text{intercept}=-6.54$), our
AEFM enumeration approach benefits greatly from parallelism. 

Not only are ACHMCs computed independently of one another, but our program
is also parallelised on the level of individual ACHMCs to improve AEFM
enumeration times. For the top 31 largest carbon ACHMCs in the HepG2
network, Figure \ref{fig03:Benchmark}d shows a median speedup of 7.9 times
when using 16 or 32 threads. 

\subsection*{Most AEFMs are source-to-sink pathways spanning dozens of reactions}

We next turned to characterising the structural properties of all carbon
and nitrogen AEFM across the five GEMs. Molecular EFMs can be divided into
pathway that span source and sink metabolites or internal loops within the
metabolic network such as pairs of reversible reactions. Both pathways
have clear biological interpretations with source-to-sink pathways
explaining the steady state flow of source metabolites on their way to
exiting the network, and internally looped pathways corresponding to
substrate remodelling/transport cycles. At an atomic level, these pathways
can be subclassified further into AEFMs that transit sets of distinct
metabolites or those that revisit the same metabolite, albeit in different
atomic positions. 
%This notion of a ``metabolite revisitation'' is similar to an isotopomer
%in stable isotope tracing~\citep{AntoniewiczMaciekR2018Agt1}. 
An example source-to-sink AEFM with this property is a glucose-derived
carbon AEFM remodelling through two turns of the TCA cycle; on the second
turn, the position of the carbon atom shifts within the TCA intermediates,
resulting in different isotopomers~\citep{DuanLikun20221tas} and therefore
a source-to-sink pathway with a metabolite revisitation under our
nomenclature. The remodelling of fructose 6-phosphate to fructose
1,6-bisphosphate via phosphofructokinase and fructose bisphosphatase
constitutes a cyclic carbon AEFM.

To obtain a broad understanding of the types of AEFMs within the five
GEMs, we categorised them into source-to-sink pathways and looped pathways
with or without metabolite revisitations in different atomic positions.
Figure \ref{fig05:Structural}a-b shows the counts for each AEFM classes
within each GEM. As expected, we observed that the majority of carbon
AEFMs were source-to-sink pathways with or without metabolite
revisitations. Those without metabolite revisitations tended to scale as
a function of network size, while those with revisitations increased
modestly. Investigation of these source-to-sink pathways confirmed that
many metabolite revisitations involved TCA cycle intermediates. Looped
AEFMs without metabolite revisitations primarily corresponded to transport
reactions cycling metabolites across compartments and cyclic pathways
involving TCA intermediates. These two observations explain, in part, why
the iAB RBC 283 network contained the fewest number of pathways with
metabolite revisitations since they lack organelles and cannot support
cyclic AEFMs associated with metabolite transport between compartments nor
mitochondrial TCA activity. There were very few internally looped AEFMs
with metabolite revisitations with no obvious trend across the five GEMs.
The greatest number of these pathways were found in the E. coli core
network with many metabolite revisitations involving carbon dioxide
derived from isocitrate and alpha-ketoglutarate re-entering the TCA cycle
via oxaloacetate. As there are fewer metabolites containing nitrogen, we
found the majority of their AEFMs corresponded to pathways without
metabolite revisitations, primarily involving amino acid remodelling or
transports between compartments. Looped pathways with metabolite
revisitations were only observed in the iIT341 and iSB619 networks through
reactions involving the production and consumption of ammonia. 

Since many looped AEFMs involve reversible reactions or metabolite
transport across compartments, we hypothesised that looped AEFMs should
involve fewer reactions than source-to-sink pathways which can span
numerous internal metabolites. We calculated the fraction of carbon and
nitrogen AEFMs with given lengths across all five GEMs and plotted their
distributions in Figure \ref{fig05:Structural}c-d. Our results show that
carbon AEFMs in the first four networks are bimodally distributed with the
majority of pathways involving a dozen or fewer reactions. While carbon
AEFMs of the HepG2 network appeared to exhibit a single Gaussian
distribution ($\mu=67.7$, $\sigma^2=16.2$), closer inspection of the
histogram revealed a small peak spanning 2-12 reactions. To determine
whether these peaks were associated with source-to-sink and looped AEFM
classes, we reclassified the carbon AEFMs spanning less than 5-12
reactions depending on the GEM (shorter pathways), and the remaining
carbon AEFMs (longer pathways). We found that longer pathways were
associated with source-to-sink pathways, with shorter pathways containing
the majority of looped pathways. In fact, 65-97\% of all longer pathways
in each GEM were source-to-sink pathways with metabolite revisitations,
with this statistic increasing to 95-100\% when including source-to-sink
pathways without metabolite revisitations. Among shorter pathways, we
observed that roughly half or more of them were looped carbon AEFMs with
or without metabolite revisitations. Analysis of the nitrogen AEFMs
revealed only a single distribution of relatively short pathways, since
there are much fewer nitrogen-containing metabolites across all GEMs.

We next investigated how atoms within the same source metabolite may
follow different paths through the network, irrespective of their AEFM
class. We first computed the number of AEFMs for each carbon atom (resp.
nitrogen atom). For each metabolite, we then computed the ratio of the
class. We then computed the number of AEFMs for each carbon atom (resp.
nitrogen atom). So for example, a ratio of 10 means that one of the carbon
atoms in the metabolite can transit the network in 10 times as many
different ways as the carbon with the fewest paths through the network.
Our results are shown in Figure \ref{fig05:Structural}e-f, where each dot
corresponds to the ratio for a given source metabolite. We found that
a surprising number of carbons and nitrogens can transit different
pathways within their respective source metabolites. Focusing on the
source metabolites with carbon and nitrogen ratios over 100, we observed
many of these corresponded to hydrophobic and positively-charged amino
acids such as phenylalanine, methionine, arginine, and histidine. Source
metabolites containing fewer carbons or nitrogens tended to exhibit ratios
closer to 1 with AEFMs involving the same metabolite subsets. Metabolites
such as glycine, for example, were equivalently metabolised across the
same carbon AEFMs. For each GEM, and not counting source metabolites with
ratios of one, we computed median carbon ratios between 2.0 and 41.8 and
nitrogen ratios between 1.3 and 11.2. Overall, these findings highlight
the varying degrees of metabolic flexibility for the atomic constituents
of source metabolites to transit the metabolic network.

\subsection*{A minority of carbon AEFMs explains the majority of source metabolite remodelling}

% Glutamine carbon 1:   6,554
% Glutamine carbon 2: 221,933
% Glutamine carbon 3: 184,823
% Glutamine carbon 4:  28,765
% Glutamine carbon 5:   6,554
% Glutamine carbon AEFM total: 448,629
% Glutamine carbons can remodel into 203 distinct metabolites aside from
% input glutamine

%Although the number of AEFMs is smaller than the number of molecular
%EFMs, there are still thousands of these pathways.

% Given the thousands of carbon AEFMs enumerated across the five GEMs, an
% important question is to determine the most important pathways assoc
% explaining the obser

% Hm. Skips too far ahead. Need to carefully lay out why AEFM analysis is
% important to begin with.

In the previous sections, we showed how to enumerate carbon and nitrogen
AEFMs in large-scale networks and analyse the structure of these pathways.
Yet, an equally important focus of our work is to determine the amount of
flux moving through AEFM. This computational question is biologically
motivated to identify the most important pathways among the thousands to
millions of AEFMs. Furthermore, AEFMs with the greatest weight, and
therefore carrying the most atomic flux, may correspond to existing
metabolic subsystems or novel metabolic pathways. In the context of the
HepG2 network, which contains a single set of steady state fluxes
estimated during the exponential growth
phase~\citep{NilssonAvlant2020Qaoa}, it is well-known that cellular
metabolism depends on the non-essential amino acid glutamine for survival
and proliferation~\citep{YeYanyan2023Gmri,YooHeeChan2020Gric}. Hence,
a AEFM analysis of glutamine carbons may reveal alternative paths of
glutamine metabolism and quantify the amount of flux moving through these
pathways.


%To identify the most important carbon AEFMs, we turned to the HepG2
%network which contains a single set of steady state fluxes estimated by
%\citeauthor{NilssonAvlant2020Qaoa}~\citep{NilssonAvlant2020Qaoa}. These
%fluxes were estimated from a parsimonious FBA model of HepG2 cells
%cultured under exponential growth with external fluxes constrained by
%time-resolved exometabolomic measurements (see Methods). 

%Our focus was to analyse this flux network using our AEFM framework to
%identify alternative pathways of glutamine metabolism since many cancer
%cells, including HepG2, are dependent on this non-essential amino acid
%for survival and proliferation~\citep{YeYanyan2023Gmri,YooHeeChan2020Gric}.


Using our full AEFM pipeline, we decomposed the set of steady state fluxes
onto the 448,629 glutamine-derived carbon AEFMs characterising the HepG2
metabolic network. Figure \ref{fig06:Hepg2}a shows the cumulative
explained mass flow of each glutamine carbon as a function of their AEFMs
ranked from greatest to smallest explained mass flow. We capped the
cumulative explained mass flow of each glutamine carbon at 99\% because we
observed that 99.5\% of the total AEFMs collectively explained the
remaining 1\% glutamine carbon flow. These results suggested that the
majority of glutamine carbon remodelling is explained by a surprisingly
small number of glutamine-derived carbon AEFMs, especially since our
Markovian flux decomposition approach assigns some weight to all network
pathways. The top 5 ranked AEFMs across each glutamine carbon (25 pathways
altogether) explained nearly half (43.9\%) of the total glutamine carbon
mass flow with the subsequent 5 next highest ranked AEFMs only explaining
an additional 8.1\% glutamine carbon mass flow. These observations were
consistent across carbon AEFMs derived from other amino acids as well as
glucose, leading us to hypothesise that steady state pathway fluxes are
centralised to a relatively small subset of AEFMs (Figure
\ref{figS06:Cumulative}). For some source metabolites, we observed
overlapping curves between carbons (e.g.\ arginine and glycine in Figure
\ref{figS06:Cumulative}), suggesting that individual carbons were
remodelled through the same AEFMs. However, corroborating results from
Figure \ref{fig05:Structural}e-f, we observed many non-overlapping curves
for the majority of amino acids since their carbons could transit the
network via different internal metabolites.

%We denote this carbon mass flow as $C_k = w_k \cdot L_k$, where $L_k$ is
%the length of EFM $k$. 

To obtain a high-level understanding of the dominant pathways of glutamine
carbon remodelling, we constructed a metabolic subnetwork from the top
5 glutamine carbon AEFMs, collapsing metabolite-atom positions onto
distinct metabolites, since none of these AEFMs involved metabolite
revisitations. Of the 435 metabolites in the HepG2 dataset, input
glutamine can remodel into 203 species localised to the cytosol or
mitochondria. We found that the top 5 glutamine-derived carbon AEFMs
consist of just 22 metabolites and 29 reactions (Figure
\ref{fig06:Hepg2}b). These AEFMs generally consisted of cyclic pathways
involving TCA intermediates. We found several source-to-sink AEFMs
involving the production of carbon dioxide via the decarboxylation of
isocitrate in the cytosol and alpha-ketoglutarate in the mitochondria.

As AEFMs are steady state pathways, they can be visualised as Sankey
diagrams to better understand the atomic fluxes through the atomic flux
subgraph. Figure \ref{fig06:Hepg2}c-e shows the top 5 AEFMs for each
glutamine carbon. Each AEFM in the diagram is shown as a flow proportional
to its AEFM weight with nodes representing the metabolite transited by the
given glutamine carbon. Although each glutamine carbon can transit
a different number of AEFMs, we found the top 5 backbone carbons for the
first three sidechain carbons were identical in AEFM pathway with
approximately the same weights.

Focusing on the first glutamine carbon in Figure \ref{fig06:Hepg2}c, we
observed that AEFMs with the greatest weight corresponded to
well-established metabolic subsystems. The first, third, and fourth AEFMs,
for example, corresponded to aspects of the malate-aspartate shuttle. The
fifth AEFM involved glutamine remodelling into alpha-ketoglutarate and
entering the canonical TCA cycle and leaving as carbon dioxide via
isocitrate. We also noticed that the second largest AEFM cycled carbon
mass flow between TCA intermediates citrate, oxaloacetate, and malate
between the cytosol and mitochondria. Interestingly, this metabolic
pathway corresponds to the recently discovered non-canonical TCA activity
described by \citeauthor{ArnoldPaigeK2022Anta}, who validated this pathway
in non-small cell lung cancer, mouse embryonic, and C2C12 mouse myoblast
cells by stable isotope tracing analysis~\citep{ArnoldPaigeK2022Anta}.
This non-canonical TCA activity was further observed in glutamine carbons
2-4 in the second AEFM (Figure \ref{fig06:Hepg2}d). We also found first,
third, fourth, and fifth AEFMs in in Figure \ref{fig06:Hepg2}d-e
corresponded to the malate-asparate shuttle. The terminal glutamine carbon
in Figure \ref{fig06:Hepg2}e contained three source-to-sink AEFMs that
either transported glutamine out of the system or decarboxylated it into
carbon dioxide via alpha-ketoglutarate.

Thus far, we have shown that our analyses revealed that the most active
glutamine-derived carbon AEFM weights corresponded to well-known metabolic
subsystems. This finding led us to ask whether we could summarise our AEFM
analysis as a simplified metabolic model describing over 40\% of the
glutamine carbon remodelling in the HepG2 flux network. We present the
resulting schematic in Figure \ref{fig06:Hepg2}f which consists
of the following three major metabolic pathways: canonical TCA cycle,
non-canonical TCA cycle, and the malate-aspartate shuttle. The overall
schematic suggests that input glutamine is primarily routed into the
mitochondria via aspartate/glutamate antiporters encoded by the
\textit{SLC25A12} and \textit{SLC25A13} annotated in the HepG2 metabolic
model. This mitochondrial glutamate is used to regenerate TCA cycle
intermediates by converting glutamate into alpha-ketoglutarate which is
subsequently consumed in the canonical TCA cycle to produce carbon dioxide
and NADH/ATP/FADH\textsubscript{2}. More importantly, a greater proportion
of glutamate appears to be recycled through the non-canonical TCA cycle,
which does not generate NADH, but actively transports existing NADH from
the cytosol into the mitochondrial. Taken together, the dominant carbon
AEFMs predicted by our model suggest HepG2 glutamine both fuels TCA cycle
activity and regulates cellular levels of NAD\textsuperscript{+}/NADH.

\section*{Discussion}


% Combination of ideas

% Limitations of molecular EFMs

% Previous work/extending Pey's ECMs to AEFMs

% AEFMs blend/merge paths in metabolite-atom graphs with the
% non-decomposable steady state property of EFMs

Since the proposal of molecular EFMs, much effort has been made to
enumerate and decompose steady state metabolic fluxes onto these pathways.
By defining EFMs with respect to individual atoms rather than molecules,
we introduced the AEFM which describes how atomic constituents are
remodelled within a metabolic flux network. As AEFMs never involve
branching pathways, they may be easier to interpret than their molecular
EFM counterparts, which may involve multiple inputs and outputs to balance
molecular stoichiometries. Using the atom mapping program RXNMapper, we
adapted our previously described CHMC method to efficiently enumerate
carbon and nitrogen AEFMs in five GEMs and characterise source metabolite
remodelling at the atomic level. By comparing the number of AEFMs across
atoms within the same source metabolite, we revealed how certain amino
acids displayed remarkable variability in the number of pathways transited
by their constituent atoms. This analysis, in particular, may be useful
for designing non-uniformly labelled stable isotope tracers to identify
fluxes within specific metabolic subsystems. Our final AEFM weight
analysis of the HepG2 network revealed that the majority of glutamine
carbons were remodelled via well-known metabolic pathways such as the
canonical TCA cycle but also novel pathways such as non-canonical TCA
activity which has not been previously described in HepG2 cells.

%As a systematic method to explain the flow of metabolites,
The problem of molecular EFM enumeration has historically limited their
applicability to small-scale networks. Many algorithmic advances have
incrementally improved molecular EFM enumeration, yet never feasibly
beyond networks with $>100$ metabolites and
reactions~\citep{KampAxelvon2006M5fa,TerzerMarco2008Lcoe,vanKlinkenJanBert2016Faet}.
Molecular EFM enumeration is generally regarded as infeasible for
large-scale networks---so much so that numerous methods have been
developed to enumerate a subset of molecular EFMs constrained by pathway
length~\citep{deFigueiredoLuisF2009Ctse}, reaction
thermodynamics~\citep{GerstlMatthiasP2015Mief}, participating
metabolites~\citep{PeyJon2014Dcoe}, participating
reactions~\citep{KaletaChristoph2009Ctwb}, gene regulatory
rules~\citep{JungreuthmayerChristian2013rSue}, metabolic
subsystems~\citep{SchusterS2002Etps,SchwartzJean-Marc2007Omfa},
extracellular flux
measurements~\citep{SoonsZitaITA2011Iomm,JungersRaphaelM.2011Fcom}, or
even random sampling~\citep{MachadoDaniel2012Rsoe}. While our HepG2
results did show the majority of source metabolite fluxes were explained
by a small number of relatively short AEFMs, contextualising the
significance of these dominant pathway fluxes is only possible when all
AEFMs are known. Furthermore, reaction thermodynamics, gene regulatory
rules, and other constraints used to identify these EFM subsets may change
across different biological conditions. By computing AEFMs instead, we
could feasibly enumerate these pathways in large-scale networks containing
hundreds of metabolites and reactions within minutes to hours on
consumer-grade computing hardware. We are hopeful that algorithmic tricks
borrowed from molecular EFM enumeration tools, such as linear pathway
compression, could further decrease memory usage and run times.

Aside from carbon and nitrogen AEFMs, our work is broadly applicable to
other types of chemical elements. One could feasibly enumerate oxygen or
hydrogen AEFMs in small-scale networks, although the significance of these
pathways is unclear for biochemical networks given the ubiquitous nature
of these atoms. We also caution against enumerating hydrogen AEFMs, in
particular, given limitations associated with RXNMapper. Unless explicitly
represented in the reaction SMILES string, RXNMapper does not assign
hydrogen atom mappings. Moreover, RXNMapper cannot map atoms in reaction
SMILES strings longer than 512 characters. Our method could equally be
applied to study sulphur metabolism to quantify methionine or cysteine
remodelling, or phosphorus metabolism to quantify nucleotide synthesis.
Outside the realm of biochemistry, our method is applicable to any
periodic table element and possibly relevant for studying general chemical
reaction networks. 

Our AEFM framework opens many new computational problems in metabolism.
While we analysed one set of fluxes for the HepG2 network, one could
envision a differential AEFM analysis of flux distributions across
distinct conditions. These differential AEFMs may reveal changes in
metabolic subsystem activity or predict new pathways of source metabolite
remodelling. For instance, analysing carbon AEFMs of C1 compounds may be
prove useful for quantifying or optimising carbon fixation pathways in
C1-utilising organisms. Another important question is whether a joint
analysis of AEFM weights could reveal broader metabolic changes within
a condition. We observed that many AEFMs across distinct source
metabolites converged on common internal metabolites. Intuitively, the
weights of AEFMs rooted on different source metabolites, but sharing
common downstream metabolite subsequences, should be related, and these
relationship may further speed AEFM enumeration or weight computations.
Finally, more work is needed to establish relationships between AEFM
decompositions and more traditional EFM decompositions of fluxes, and to
explore the possibility that feasible AEFM calculations might somehow aid
in molecular EFM calculations.

\section*{Materials and Methods}

%\subsection*{Reagents and Tools table}

%See tables/reagents-and-tools-table.xls.

\subsection*{Methods and Protocols}

The following subsections describe in greater detail the datasets and
pipeline steps referenced in Figure~\ref{fig02:Pipeline}.

\subsubsection*{Step 1a: Datasets}

Five highly-curated GEMs were chosen in this study to reflect a diversity
of organisms, metabolites and biochemical reactions. Four networks were
obtained from the BiGG UCSD database~\citep{KingZacharyA2016BMAp} in
ascending order of network size. Datasets were downloaded as SBML (.xml)
files, and the corresponding stoichiometry matrices, metabolites and
reactions were extracted using the Julia SBML.jl package. The HepG2 liver
cancer cell line dataset was chosen for its single set of estimated steady
state fluxes estimated from a FBA with external fluxes constrained by
time-resolved exometabolomic measurements of glucose, pyruvate, lactate,
and various amino acids. These cells were grown in DMEM with 10\% fetal
calf serum, 22\SI{}{\milli\molar} glucose, and 1.8\SI{}{\milli\molar}
glutamine~\citep{NilssonAvlant2020Qaoa}. The steady state fluxes were
computed from a piecewise FBA model provided by the authors
(figure3\_metabolitebalances.m; MATLAB R2018a)
(\citealt{NilssonAvlant2020Qaoa}; Data ref:
\citealt{NilssonAvlant2020QaoaDATASET}). Across all models, we modelled
all reversible reactions as net forward reactions to simplify the number
of molecular and AEFMs as the majority of experimental and computational
methods can only estimate net fluxes. We assumed the reaction
stoichiometries of the BiGG datasets represented the canonical net
directions of all reversible reaction. For the HepG2 dataset, the net
directions were chosen based on the net fluxes estimated from the FBA
model.

\subsubsection*{Step 1b: Extracting metabolites SMILES strings}

SMILES structures were not encoded in the BiGG SBML files nor the HepG2
dataset. Thus, we manually extracted SMILES strings corresponding to each
metabolite across all GEMs. Isomeric SMILES strings were taken from
PubChem~\citep{KimSunghwan2023P2u}, the Human Metabolome
Database~\citep{WishartDavidS2022H5tH}, or MetaNetX
database~\citep{MorettiSebastien2021Munf}. Metabolites shared across GEMs
were assigned the same SMILES structure in addition to identical
metabolites stratified across distinct subcellular compartments. Generic
structures such as R-groups were not modelled in this study and were
replaced with a single bonded hydrogen under the assumption this
modification would not alter atom mappings, especially since RXNMapper
does not map hydrogen atoms unless they are explicitly represented in the
SMILES strings. SMILES strings were also not assigned to metabolites with
ambiguous structures nor pseudometabolites with no known chemical
structure. All pseudometabolites and reactions involving pseudometabolites
were removed from the stoichiometry matrix (Table 2).

\subsubsection*{Step 1c: Constructing reaction SMILES strings}

RXNMapper performs atom mapping on a string that encodes the SMILES
structures of all substrates and products for a given reaction. These
reaction SMILES strings are constructed by concatenating SMILES strings of
individual substrates and products. Multiple substrates or products are
separated by the symbol `.', while the reaction arrow delimiting
substrates and products is denoted by the symbol `$>>$'. The substrate and
product SMILES strings are concatenated in order of metabolite appearance
within the stoichiometry matrix, and multiple stoichiometric copies are
consecutively repeated in the reaction SMILES strings. Non-integer
stoichiometric coefficients are therefore not allowed and these
pseudoreactions were removed from the stoichiometry matrix (see Table 2).
To maintain steady state flux in the HepG2 network, flux along
pseudoreactions were replaced with unimolecular source and sink fluxes
consuming each pseudoreaction substrate and producing each pseudoreaction
product.

\subsubsection*{Step 1d: Pre-processing genome-scale metabolic models}

Several pre-processing steps were required to compute AEFMs from the GEMs.
The stoichiometric coefficients of external reactions were set to one and
their fluxes rescaled to carry equivalent units of source and sink
metabolites. All reactions containing pseudometabolites with no known
SMILES strings were removed from the stoichiometry matrix and replaced
with unimolecular source and sink reactions carrying equivalent units of
flux. The same procedure was also applied to reactions containing
non-integer stoichiometries. These pseudometabolites were next removed
from the stoichiometry matrix rows, and unimolecular reactions involving
the same external or internal metabolites were aggregated. Finally, we had
to remove a small number of reactions whose reaction SMILES strings
exceeded the 512 token limit of RXNMapper. These reactions were likewise
converted into unimolecular source and sink reactions to maintain steady
state flux across the network (see Table 2). All pre-processing steps
described above are wrapped into convenience functions within our Julia
package MarkovWeightedEFMs.jl.

\subsubsection*{Step 1e: running RXNMapper to map atomic movements across
reactions}

After GEM pre-processing, RXNMapper was used to predict the atom mappings
across reaction SMILES strings. The output of RXNMapper is a mapped
reaction SMILES string with numeric indices indicating the atomic
movements from substrate to product side of each reaction SMILES string.
We manually checked several reactions to ensure the assigned atom
mappings were correct. Crucially, we found RXNMapper to fail on nearly all
classes of ATP-independent transaminase reactions. This problem has
previously been reported in an atom mapping benchmark study prior to the
publication of RXNMapper~\citep{PreciatGonzalezGermanA.2017Ceoa}. These
transaminase reactions were identified in all metabolic networks and
manually atom mapped based on the KEGG reaction mappings (Figure
\ref{figS07:transaminase}). Additionally, the reactions converting citrate
to isocitrate (aconitate hydratase) and alpha-ketoglutarate to
succinyl-CoA required manual curation.

We then store these curated atom mapping predictions in a GEM and atom
type-specific atom mapping dictionary. Each key corresponds to
a substrate-atom uniquely defined by its (i) metabolite index within the
stoichiometry matrix, (ii) ordinal occurrence of that atom type within the
substrate SMILES string, (iii) substrate stoichiometric copy, and (iv)
reaction index within the stoichiometry matrix. The values of this
dictionary are the corresponding product-atom positions uniquely defined
by the (i) metabolite index within the stoichiometry matrix and (ii)
ordinal occurrence of that atom type in the product SMILES string.

\subsubsection*{Step 2: Constructing atomic flux subgraphs}

%\subsubsection*{Step 2a: mapping atomic movements across reactions}

%\subsubsection*{Step 2b: constructing atomic flux subgraphs}

%\hl{

From the atom-mapping dictionary, we then construct atomic flux subgraphs
rooted on each source metabolite-atom position. Starting with the source
metabolite-atom position, we search the atom-mapping dictionary to
determine the participating reactions and product metabolite-atom
positions. For reactions involving a single stoichiometric copy of each
substrate and product, the atom of interest will always transition from
a unique substrate metabolite-atom position to product metabolite-atom
position. If this atom is present in a reaction with multiple substrate
stoichiometric copies, we assume there is an equal probability of that
atom occurring in either substrate stoichiometric copy. Depending on the
substrate stoichiometric position, the atom may (i) transition to
different product metabolite-atom positions, (ii) the same product
metabolite but at different atomic positions, or (iii) the same product
metabolite-atom position. In all cases, we set the atomic flux of the atom
transitioning from either stoichiometric copy equal to the reaction fluxes
to maintain atomic mass conservation (Figure {\ref{figS01:Rules1}}). The
resulting atomic flux graph is completed by connecting any sink fluxes to
the root of the tree initialised on the source metabolite-atom position.
This final procedure ensures a closed-loop subgraph maintaining atomic
mass conservation.

%}

\subsubsection*{Step 3: Building the ACHMC model and enumerating AEFMs}

%\hl{
The atomic flux subgraphs described in the previous step satisfy all
conditions required to construct the corresponding ACHMC. Specifically,
these conditions are that the input flux network is a closed-loop graph
that is strongly connected with conserved edge weights such that the sum
of weights entering and exiting each node are
balanced~\citep{ChitpinJustinG2023AMct}.

Rather than explicitly constructing an atomic flux subgraph for each
source metabolite-atom position and the corresponding ACHMC, we choose to
construct the ACHMCs directly from the metabolic flux network and the atom
mapping dictionary. Starting at an initial metabolite-atom position, we
apply a depth-first search algorithm (Algorithm~{\ref{alg01:trie}}) to
recursively enumerate sequences of transited metabolite-atom positions
stemming from the initial source metabolite atom. Following
\citeauthor{ChitpinJustinG2023AMct}, we denote state $S$ in the ACHMC to
represent the simple path of metabolite-atom transitions in the trie $s_0,
s_1, \ldots, s_m$. Linear AEFMs (simple paths) are enumerated by finding
all states $S$ (sequences of metabolite-atom positions) that can no longer
transition to a longer simple path (leaves within the trie).
Internally-looped AEFMs (simple cycles) are enumerated by finding all
prefixes that are capable of transitioning to a shorter simple path within
the trie. 

%}

\begin{algorithm}
\caption{Depth-first construction of the ACHMC trie}\label{alg01:trie}
\hspace*{\algorithmicindent} \textbf{Input:} Input state $s_{im}$
denoting an atom in position $i$ of metabolite $m$, dictionary $D$ \\
%\hspace*{\algorithmicindent} mapping $\{a_{im}\}_r$ to $\{a_{jn}\}_r$ in reaction $r$, stoichiometry matrix $S = [s_{mn}]$, steady state\\
\hspace*{\algorithmicindent} mapping $\{s_{im}\}_r$ to $\{s_{jn}\}_r$ in reaction $r$, stoichiometry matrix $S$, steady state reaction\\
\hspace*{\algorithmicindent} flux vector $v$\\
\hspace*{\algorithmicindent} \textbf{Output:} Dictionary $T$ encoding the ACHMC trie
\begin{algorithmic}[1]
\Function{traverse\_trie}{$s_{im}$, $D$, $S$, $v$}
    \State $V \gets \{s_{im}\} \text{ Array of visited metabolite-atom positions}$
    \State $W \gets V[\text{end}] \text{}$
    \State $T \gets \{\}\ \text{ where T[\text{Prefix}]} \gets \{\text{Children}, \text{ Children values}, \text{ Parent}, \text{ Parent values}\}$
    \State $R \gets \text{reaction indices in } S \text{ where } W \text{ participates as a substrate}$
    \State $s_{out} \gets \{D[W]\}_{r \in R}$
    \State $s_{down} \gets s_{out} \notin V$ 
    \State $s_{up} \gets s_{out} \in V$
    \State $\text{Add } \{s_{down}, v[r\in R], s_{up}, v[r\notin R]\} \text{ to } T[V]$
    \For{$s \in s_{down}$}
        \State $\textsc{traverse\_trie}(\{V, s\}, D, S, v)$
    \EndFor
\EndFunction
\end{algorithmic}
\end{algorithm}

\subsubsection*{Step 4: Computing AEFM weights in the HepG2 network}

%\hl{
To compute AEFM weights, we encode the ACHMC trie into a sparse transition
probability matrix $T$. Each element $T_{S,S^\prime}$ corresponds to the
probability of an atom in state $S \equiv s_0, s_1, \ldots, s_m$
transitioning to either state $S^\prime \equiv s_0, s_1, \ldots, s_m,
s_{m+1}$ or state $S^\prime \equiv s_0, s_1, \ldots, s_i$ if $s_{m+1}$ is
among $s_0, s_1, \ldots, s_m$. If no such state $s_{m+1}$ exists, we set
$T_{S,S^0} = 1$ to ensure the ACHMC is strongly connected where $S^0
\equiv s_0$. The transition probabilities are taken from the atomic flux
values stored within the trie and normalised such that
$\sum_{S^{\prime}\notin S} T_{S,S^{\prime}} = 1$.

The AEFM probabilities are computed by steady state analysis of the ACHMC
transition probability matrix~\citep{ChitpinJustinG2023AMct}. For AEFM
$k$, we identify all transitions in the ACHMC that close the simple cycle
$E_k$. For linear AEFMs, $E_k$ always corresponds to one transition;
however, internally looped AEFMs may contain multiple transitions within
$E_k$. The steady state probability of transiting AEFM $k$ is defined as
the sum of steady state probabilities of those transitions or:
%}
%
\begin{equation}
    P_k = \sum_{(S,S^\prime) \in E_k} \pi_S \cdot T_{S, S^\prime}
\end{equation}
%
%\hl{
These AEFM probabilities are scaled to weights under the conservation
constraint that the atomic flux $v^\prime_0$ entering a given source
metabolite is equal to the observed (unimolecular) reaction flux $v_0$ of
the underlying metabolic flux network. Thus for each ACHMC, we identify
all $l$ linear AEFMs and compute the ACHMC-specific scaling constant
$\alpha$ satisfying $v_0 = \sum_{k \in l}P_k \cdot \alpha$. 
%}

\subsubsection*{Benchmarking the enumeration of atomic-versus-standard EFMs}

We benchmarked our Julia package MarkovWeightedEFMs.jl version 2.0
running Julia version 1.10.2 to enumerate
AEFMs versus the molecular EFM enumeration program FluxModeCalculator
(\cite{vanKlinkenJanBert2016Faet}) running MATLAB version R2018a with
default parameters including network compression. We chose
FluxModeCalculator because it is the state-of-the-art program for
enumerating molecular EFMs. Benchmarks were performed on a 64-bit Linux
operating system with 64GB of memory and an AMD 5950X CPU with 16 cores/32
threads. All benchmarks reported the wall time to run the AEFM or
molecular EFM enumeration functions and did not include startup times to
open Julia/MATLAB, precompile Julia functions, or load packages and input
data. A benchmark time of `did not finish' (DNF) was assigned if either
program could not complete (A)EFM enumeration within 7 days. In practice,
only FluxModeCalculator failed to enumerate the molecular EFMs in three of
the five datasets.


\section*{Data availability}

Our software package MarkovWeightedEFMs.jl is available at
\url{https://github.com/jchitpin/MarkovWeightedEFMs.jl}. All source codes
to reproduce this study are available at
\url{https://www.github.com/jchitpin/reproduce-efm-paper-2024}.

\section*{Acknowledgements}


%
This research was supported by the Natural Sciences and Engineering
Research Council of Canada (NSERC), Discovery grant RGPIN/06604-2019 and
a Digital Research Alliance of Canada (DRAC) Resources for Research Groups
grant to TJP (RRG-TPERKINS). J.G.C. was supported by an NSERC CREATE
Matrix Metabolomics Scholarship, NSERC Alexander Graham Bell Canada
Graduate Scholarship, and Ontario Graduate Scholarship.
%\hl{
The authors also thank the anonymous reviewers for their helpful comments.
%}

\section*{Author contributions}

Justin G. Chitpin: Conceptualisation; data curation; methodology;
software; formal analysis; investigation; visualisation; writing
- original draft; writing - review and editing.\\

\noindent Theodore J. Perkins: Conceptualisation; methodology;
supervision; funding acquisition; investigation; project administration;
writing - review and editing.

\section*{Disclosure and competing interests statement}

The authors declare no competing interests.

\bibliography{biblio-aefms}


\newpage
\pagenumbering{gobble}

\begin{figure}
\includegraphics[width=\textwidth]{Fig/figure-network-notions-efms-aefms.pdf}
\caption{Various notions of metabolic networks with respect to an (a)
example metabolic flux network, (b) metabolic graph, and (c) atomic
transition/flux graph. The corresponding elementary flux modes are shown
for the (d) metabolic flux network and (e) atomic transition/flux
graph. \label{fig01:Notions}}
\end{figure}

\begin{figure}
\includegraphics[width=\textwidth]{Fig/figure-pipeline.pdf}
\caption{Atomic cycle-history Markov chain (ACHMC) pipeline. (a) An
example molecular EFM in glycolysis. (b) An example AEFM tracing the flow
of a given glucose carbon in glycolysis. (c) The major steps within our
pipeline to enumerate AEFMs and compute their weights. \label{fig02:Pipeline}}
\end{figure}

\begin{figure}
\includegraphics[width=\textwidth]{Fig/figure-benchmarking.pdf}
\caption{AEFM enumeration is computationally tractable
compared to molecular EFMs in five GEMs. (a) Number of atomic and
molecular EFMs in the five GEMs. (b) The run time of our
MarkovWeightedEFMs.jl method versus FluxModeCalculator running in serial.
(c) Serial run time scaling of MarkovWeightedEFMs.jl as a function of the
number of AEFMs. (d) Speedup of selected ACHMCs as a function of
additional threads. \label{fig03:Benchmark}}
\end{figure}

\begin{figure}
\centerline{\includegraphics[height=8in]{Fig/figure-structural.pdf}}
%\centerline{\includegraphics[width=\textwidth]{Fig/figure-structural.pdf}}
\caption{Structural analysis of AEFMs in five genome-scale
metabolic models. (a-b) Classes of carbon and nitrogen EFMs. (c-d)
Fraction of carbon and nitrogen AEFM lengths. (e-f) Ratio of the greatest
and smallest number of AEFMs between carbons and nitrogens within the same
source metabolite. Representative metabolites from each dataset are
labelled; labels with asterisks indicate that metabolite is also present
in other networks. \label{fig05:Structural}}
\end{figure}

\begin{figure}
\includegraphics[width=\textwidth]{Fig/figure-hepg2-weights.pdf}
\caption{Carbon AEFM weight analysis of the HepG2 metabolic
flux network. (a) A small number of highly active AEFMs explain the
majority of glutamine carbon flow. (b) Subnetwork of the top five AEFMs across
each glutamine carbon collapsed onto metabolites. (c-e) Sankey diagrams of
the top five glutamine carbon AEFMs. (f) Schematic mapping AEFMs onto
traditionally defined metabolic pathways/subsystems. All metabolite names are
labelled following the BiGG database nomenclature. \label{fig06:Hepg2}}
\end{figure}

%\section*{Tables and their legends}

% Table 1
\begin{table}
\caption{Summary statistics for the five metabolic networks in this study.}
~\\
\begin{tabular}{l|ccccc}
& \multicolumn{5}{c}{{Carbon metabolism}} \\
{GEM (\# mets $|$ \# rxns)} & {Sources} & {Sinks} & {States} & {Transitions} & {AEFMs} \\ \hline
E. coli core (71 $|$ 144) & 184 & 255 & 14659 & 24355 & 6519 \\
iAB\_RBC\_283	 (274 $|$ 465) & 954 & 1166 & 37122 & 51220 & 7958 \\
iIT341 (433 $|$ 759) & 1486 & 2134 & 86342 & 120545 & 30633 \\
iSB619 (547 $|$ 974) & 2188 & 2986 & 169877 & 246712 & 59444 \\
HepG2 (435 $|$ 545) & 376 & 939 & 14159860 & 18152623 & 2180613\\
~\\
& \multicolumn{5}{c}{{Nitrogen metabolism}} \\
{GEM (\# mets $|$ \# rxns)} & {Sources} & {Sinks} & {States} & {Transitions} & {AEFMs} \\ \hline
E. coli core (71 $|$ 144) & 30 & 40 & 246 & 448 & 204 \\
iAB\_RBC\_283	 (274 $|$ 465) & 161 & 214 & 6409 & 9275 & 2091\\
iIT341 (433 $|$ 759) & 314 & 447 & 27440 & 39761 & 10440 \\
iSB619 (547 $|$ 974) & 353 & 516 & 36661 & 54165 & 14214 \\
HepG2 (435 $|$ 545) & 72 & 190 & 19137 & 27860 & 5105 \\
\end{tabular}
\end{table}

%Table 2:
\begin{table}
\caption{Number of metabolites and reactions removed from each network during pre-processing.}
\begin{tabular}{l|ccccc}
 & {E. coli core} & {iAB\_RBC\_283} & {iIT341} & {iSB619} & {HepG2} \\ \hline
 ~\vspace{-1ex}\\
{\small Rxns with non-integer stoichiometries} & 2 & 1 & 3 & 34 & 8 \\
{\small Pseudometabolites with no known structures} & 1 & 68 & 52 & 108 & 33 \\
{\small Rxns with pseudometabolites} & 2 & 142 & 93 & 204 & 46 \\
{\small Rxns exceeding RXNMapper token limit} & 0 & 1 & 4 & 14 & 3 
\end{tabular}
\end{table}

%\section*{Expanded View Figure legends}
\setcounter{figure}{0}
\renewcommand{\thefigure}{S\arabic{figure}}

% Expanded View 1
\begin{figure}
%\centerline{\includegraphics[height=6.25in]{Expanded/figure-atomic-chmc-cartoon-rules-top.pdf}}
\centerline{\includegraphics[width=\textwidth]{Expanded/figure-atomic-chmc-cartoon-rules-top.pdf}}
\caption{Rules for assigning atomic fluxes from
atom-mapped reactions with biological examples. Schematic node colours
represent ground truth atom mappings across substrate(s) and product(s).
Node labels correspond to distinct atomic indices of substrate and product
stoichiometric copies.
(a) For reactions with single stoichiometric copies of substrates and
products, each atom within a given substrate moves to a single atomic
position within a product with atomic flux equal to the reaction flux.
(b-d) For reactions with $s$ stoichiometric copies of a given substrate,
there are $s$ copies of each metabolite-atom state on the substrate and
product side of the reaction. The probability of any atom within
a substrate stoichiometric copy moving to a product metabolite-atom state
remains equal to the reaction flux. (b) The stoichiometric copies of the
substrate fuse together into a single product. (c) The stoichiometric
copies of the substrate form distinct products. (d) The $s$ stoichiometric
copies of the substrate form $s$ copies of products with identical atomic
backbones. \label{figS01:Rules1}}
\end{figure}

% Expanded View 2
\begin{figure}
%\centerline{\includegraphics[height=6.25in]{Expanded/figure-atomic-chmc-cartoon-rules-bottom-inkscape.pdf}}
\centerline{\includegraphics[width=\textwidth]{Expanded/figure-atomic-chmc-cartoon-rules-bottom-inkscape.pdf}}
\caption{Biological examples for rules in Figure~\ref{figS01:Rules1}.
(a) Biological example of rule a where the atoms in each
substrate map to a single position on the corresponding product. (b)
Biological example of rule b where two stoichiometric copies of
glutathione fuse together to form a single stoichiometric copy of
glutathione disulfide. (c) Biological example of rule c where two
stoichiometric copies of acetyl-CoA undergo a reaction to produce
a stoichiometric copy of two distinct products. \label{figS02:Rules2}}
\end{figure}

% Expanded View 3
\begin{figure}
\centerline{\includegraphics[height=7.5in]{Expanded/figure-gem-hairballs.pdf}}
\caption{Network visualisation of the five GEMs after
pre-processing. Metabolites (non-grey nodes) are connected to reaction
entities (grey nodes) with node sizes proportional to their degree. Graphs
were constructed in Gephi version 0.10 using the Yifan-Hu proportional
network layout with default parameters. (a) E. coli core. (b) iAB RBC 283.
(c) iIT341. (d) iSB619. (e) HepG2. \label{figS03:Hairballs}}
\end{figure}

% Expanded View 4
\begin{figure}
\centerline{\includegraphics[width=3.5in]{Expanded/figure-benchmarking-fluxmodecalculator-parallel.pdf}}
\caption{FluxModeCalculator scaling across multiple threads. \label{figS04:Scaling}}
\end{figure}

% Expanded View 5
\begin{figure}
\centerline{\includegraphics[width=\textwidth]{Expanded/figure-atomic-state-space.pdf}}
\caption{Most ACHMCs traverse less than 10\% of the
total metabolite-carbon/nitrogen state space. Each point corresponds to an
ACHMC rooted on a given source metabolite carbon or nitrogen. \label{figS05:StateSpace}}
\end{figure}

% Expanded View 6
\begin{figure}
\centerline{\includegraphics[width=\textwidth]{Expanded/figure-cumulative-combined-legend.pdf}}
\caption{Cumulative explained carbon mass flow of input glucose and amino
acid carbons from the HepG2 flux network. The cumulative explained mass
flow is truncated at 99\% to avoid displaying AEFMs that carry little
atomic mass flow. \label{figS06:Cumulative}}
\end{figure}

% Expanded View 7
\begin{figure}
%\centerline{\includegraphics[height=8in]{Expanded/figure-transaminase-examples.pdf}}
\centerline{\includegraphics[width=\textwidth]{Expanded/figure-transaminase-examples.pdf}}
\caption{Examples of incorrectly atom-mapped transaminase
reactions present in one or more of the five networks. (a) Phenylalanine
transaminase. (b) Isoleucine transaminase. (c) Tyrosine transaminase. (d)
Aspartate transaminase.\label{figS07:transaminase}}
\end{figure}

\end{document}

